\documentclass{beamer}
\usepackage[utf8]{inputenc}
\usepackage[spanish]{babel}
\usepackage{amsmath, amssymb, booktabs, array, xcolor, mathtools, graphicx, tikz}
\usetikzlibrary{arrows.meta, positioning}
\usepackage{pgfplots}
\pgfplotsset{compat=1.18}

\usetheme{Madrid}
\usecolortheme{seahorse}

\setbeamertemplate{navigation symbols}{}
\graphicspath{{images/}{./}}

\title{SuperDFS Heur\'istica}
\subtitle{Exploraci\'on en Profundidad Guiada por Pila e Ignorados}
\author{Dario Turco}
\date{}

% ---- colors ----
\colorlet{cChosen}{green!55!black}
\colorlet{cStack}{orange!85!black}
\colorlet{cIgnore}{blue!70!black}
\colorlet{cError}{red!70!black}
% ---- heatmap bucket colors (transitions) ----
\definecolor{cB1}{RGB}{255,235,0}    % <=100: yellow
\definecolor{cB2}{RGB}{255,130,0}    % 100-999: orange
\definecolor{cB3}{RGB}{200,40,80}    % 1k-5k: red
\definecolor{cB4}{RGB}{110,0,180}    % 5k-15k: violet

\newcommand{\imgpending}[1]{%
  \begin{tikzpicture}
    \draw[dashed, gray!60, rounded corners] (0,0) rectangle (9,4.5);
    \node[gray!70, font=\large] at (4.5,2.25) {\textit{Pendiente: #1}};
  \end{tikzpicture}%
}

\tikzset{
  st/.style     = {circle, draw=black, thick, minimum size=7mm, font=\small},
  errst/.style  = {circle, draw=cError, thick, minimum size=7mm, font=\small, text=cError},
  tChosen/.style = {-{Stealth[length=2.5mm]}, thick, cChosen},
  tStack/.style  = {-{Stealth[length=2.5mm]}, thick, cStack},
  tIgnore/.style = {-{Stealth[length=2.5mm]}, thick, cIgnore},
  tError/.style  = {-{Stealth[length=2.5mm]}, thick, cError},
  tPend/.style   = {-{Stealth[length=2.5mm]}, thick, gray!50},
  lbl/.style    = {font=\footnotesize},
}

\begin{document}

%--- T\'itulo ---
\begin{frame}
\titlepage
\end{frame}

%=======================================================================
% RESUMEN
%=======================================================================
\begin{frame}{Funcionamiento de la Heur\'istica}

\begin{alertblock}{Premisa}
La instancia es \textbf{resoluble} (existe un director). El estado actual $s$ \textbf{pertenece al director}.
\end{alertblock}

\medskip
DFS guiado por dos estructuras auxiliares:

\bigskip
\begin{columns}[t]
\column{0.48\textwidth}
\textbf{Estructuras:}
\begin{itemize}
  \item \textcolor{cStack}{\textbf{Pila} (\textsc{stack})} -- no-controlables diferidas
  \item \textcolor{cIgnore}{\textbf{Ignorados} (\textsc{ignore})} -- controlables no elegidas
\end{itemize}

\column{0.48\textwidth}
\textbf{En cada estado $s$:}
\begin{enumerate}
  \item Si hay no-controlables desde $s$:\\
    elegir la primera, apilar el resto
  \item Si solo hay controlables desde $s$:\\
    \emph{picker} de familia, resto $\to$ ignorados
  \item Sin transiciones desde $s$:\\
    pop \textcolor{cStack}{\textsc{stack}} $\to$ pop \textcolor{cIgnore}{\textsc{ignore}}
\end{enumerate}
\end{columns}

\end{frame}

%=======================================================================
% CASO 1 -- una controlable
%=======================================================================
\begin{frame}{Caso 1 -- \'Unica transici\'on controlable}
\begin{columns}[c]
\column{0.42\textwidth}
\begin{center}
\begin{tikzpicture}[node distance=2.2cm]
  \node[st, fill=gray!12] (s) {$s$};
  \node[st, right of=s]   (t) {$t$};
  \draw[tChosen] (s) -- node[above, lbl]{$a$} (t);
  \node[below=0.05cm of s, cChosen, font=\tiny] {actual};
\end{tikzpicture}
\end{center}

\column{0.54\textwidth}
\textbf{Desde $s$}: solo una transici\'on controlable $a$.

\medskip
\begin{itemize}
  \item[\textcolor{cChosen}{$\blacktriangleright$}] Elegir $a \to t$
  \item \textsc{stack}: sin cambios
  \item \textsc{ignore}: sin cambios
  \item Siguiente estado: $t$
\end{itemize}

\medskip
{\small No hay decisi\'on que tomar; \'unica opci\'on v\'alida.}
\end{columns}
\end{frame}

%=======================================================================
% CASO 2 -- una no-controlable
%=======================================================================
\begin{frame}{Caso 2 -- \'Unica transici\'on no-controlable}
\begin{columns}[c]
\column{0.42\textwidth}
\begin{center}
\begin{tikzpicture}[node distance=2.2cm]
  \node[st, fill=gray!12] (s) {$s$};
  \node[st, right of=s]   (t) {$t$};
  \draw[tChosen] (s) -- node[above, lbl]{${\sim}a$} (t);
  \node[below=0.05cm of s, cChosen, font=\tiny] {actual};
\end{tikzpicture}
\end{center}

\column{0.54\textwidth}
\textbf{Desde $s$}: solo una transici\'on no-controlable ${\sim}a$.

\medskip
\begin{itemize}
  \item[\textcolor{cChosen}{$\blacktriangleright$}] Elegir ${\sim}a \to t$ (forzosa)
  \item \textsc{stack}: sin cambios
  \item \textsc{ignore}: sin cambios
  \item Siguiente estado: $t$
\end{itemize}

\medskip
{\small El director \emph{siempre} debe explorar no-controlables; no puede bloquearlas.}
\end{columns}
\end{frame}

%=======================================================================
% CASO 3 -- m\'ultiples no-controlables
%=======================================================================
\begin{frame}{Caso 3 -- M\'ultiples transiciones no-controlables}
\begin{columns}[c]
\column{0.44\textwidth}
\begin{center}
\begin{tikzpicture}[node distance=1.9cm]
  \node[st, fill=gray!12] (s)  {$s$};
  \node[st, above right=0.8cm and 1.8cm of s] (t1) {$t_1$};
  \node[st, right of=s]                       (t2) {$t_2$};
  \node[st, below right=0.8cm and 1.8cm of s] (t3) {$t_3$};
  \draw[tChosen] (s) -- node[above, lbl, sloped]{${\sim}a$} (t1);
  \draw[tStack]  (s) -- node[above, lbl]{${\sim}b$}         (t2);
  \draw[tStack]  (s) -- node[below, lbl, sloped]{${\sim}c$} (t3);
  \node[below=0.05cm of s, cChosen, font=\tiny] {actual};
\end{tikzpicture}
\end{center}

\column{0.52\textwidth}
\textbf{Desde $s$}: tres no-controlables.

\medskip
\begin{itemize}
  \item[\textcolor{cChosen}{$\blacktriangleright$}] Elegir ${\sim}a \to t_1$ (primera)
  \item[\textcolor{cStack}{$\blacktriangleright$}] Apilar ${\sim}b \to t_2$
  \item[\textcolor{cStack}{$\blacktriangleright$}] Apilar ${\sim}c \to t_3$
  \item \textsc{ignore}: sin cambios
  \item Siguiente estado: $t_1$
\end{itemize}

\medskip
\begin{block}{Observaci\'on clave}
Como $s$ pertenece al director, \emph{todas} sus no-controlables tambi\'en pertenecen al director. No importa cu\'al se elige primero: el orden es irrelevante.
\end{block}
\end{columns}
\end{frame}

%=======================================================================
% CASO 4 -- m\'ultiples controlables
%=======================================================================
\begin{frame}{Caso 4 -- M\'ultiples transiciones controlables}
\begin{columns}[c]
\column{0.44\textwidth}
\begin{center}
\begin{tikzpicture}[node distance=1.9cm]
  \node[st, fill=gray!12] (s)  {$s$};
  \node[st, above right=0.8cm and 1.8cm of s] (t1) {$t_1$};
  \node[st, right of=s]                       (t2) {$t_2$};
  \node[st, below right=0.8cm and 1.8cm of s] (t3) {$t_3$};
  \draw[tChosen] (s) -- node[above, lbl, sloped]{$a$} (t1);
  \draw[tIgnore] (s) -- node[above, lbl]{$b$}         (t2);
  \draw[tIgnore] (s) -- node[below, lbl, sloped]{$c$} (t3);
  \node[below=0.05cm of s, cChosen, font=\tiny] {actual};
\end{tikzpicture}
\end{center}

\column{0.52\textwidth}
\textbf{Desde $s$}: tres controlables.

\medskip
\begin{itemize}
  \item[\textcolor{cChosen}{$\blacktriangleright$}] Elegir $a \to t_1$ (\emph{picker} de familia)
  \item[\textcolor{cIgnore}{$\blacktriangleright$}] Ignorar $b \to t_2$
  \item[\textcolor{cIgnore}{$\blacktriangleright$}] Ignorar $c \to t_3$
  \item \textsc{stack}: sin cambios
  \item Siguiente estado: $t_1$
\end{itemize}

\medskip
\begin{block}{Heur\'istica de segundo nivel}
O sea, tenemos otra heur\'istica para decidir qu\'e expandir (solo para controlables).
\end{block}
\end{columns}
\end{frame}

%=======================================================================
% CASO 5 -- mezcla no-controlables y controlables
%=======================================================================
\begin{frame}{Caso 5a -- Mezcla: picker elige no-controlable (\texttt{null})}
\begin{columns}[c]
\column{0.44\textwidth}
\begin{center}
\begin{tikzpicture}[node distance=1.9cm]
  \node[st, fill=gray!12] (s)  {$s$};
  \node[st, above right=0.8cm and 1.8cm of s] (t1) {$t_1$};
  \node[st, right of=s]                       (t2) {$t_2$};
  \node[st, below right=0.8cm and 1.8cm of s] (t3) {$t_3$};
  \draw[tChosen] (s) -- node[above, lbl, sloped]{${\sim}a$} (t1);
  \draw[tStack]  (s) -- node[above, lbl]{${\sim}b$}         (t2);
  \draw[tPend]   (s) -- node[below, lbl, sloped]{$c$}       (t3);
  \node[below=0.05cm of s, cChosen, font=\tiny] {actual};
  \node[right=0cm of t3, gray!60, font=\tiny] {(frontera)};
\end{tikzpicture}
\end{center}

\column{0.52\textwidth}
\textbf{Picker devuelve \texttt{null}}: no quiere elegir controlable.

\medskip
\begin{itemize}
  \item[\textcolor{cChosen}{$\blacktriangleright$}] Elegir ${\sim}a \to t_1$
  \item[\textcolor{cStack}{$\blacktriangleright$}] Apilar ${\sim}b \to t_2$
  \item[\textcolor{gray!70}{$-$}] $c \to t_3$ queda en la \emph{frontera}
  \item \textsc{ignore}: sin cambios
\end{itemize}

\medskip
{\small El director restringe todas las controlables; solo habilita las no-controlables.}
\end{columns}
\end{frame}

\begin{frame}{Caso 5b -- Mezcla: picker elige controlable (non-\texttt{null})}
\begin{columns}[c]
\column{0.44\textwidth}
\begin{center}
\begin{tikzpicture}[node distance=1.9cm]
  \node[st, fill=gray!12] (s)  {$s$};
  \node[st, above right=0.8cm and 1.8cm of s] (t1) {$t_1$};
  \node[st, right of=s]                       (t2) {$t_2$};
  \node[st, below right=0.8cm and 1.8cm of s] (t3) {$t_3$};
  \draw[tChosen] (s) -- node[above, lbl, sloped]{$c$}       (t1);
  \draw[tStack]  (s) -- node[above, lbl]{${\sim}a$}         (t2);
  \draw[tStack]  (s) -- node[below, lbl, sloped]{${\sim}b$} (t3);
  \node[below=0.05cm of s, cChosen, font=\tiny] {actual};
\end{tikzpicture}
\end{center}

\column{0.52\textwidth}
\textbf{Picker devuelve} $c$: elige la controlable.

\medskip
\begin{itemize}
  \item[\textcolor{cChosen}{$\blacktriangleright$}] Elegir $c \to t_1$ (\emph{picker} de familia)
  \item[\textcolor{cStack}{$\blacktriangleright$}] Apilar ${\sim}a \to t_2$
  \item[\textcolor{cStack}{$\blacktriangleright$}] Apilar ${\sim}b \to t_3$
  \item \textsc{ignore}: otras controlables (si hubiera)
\end{itemize}

\medskip
{\small Las no-controlables van al \textsc{stack} (no al \textsc{ignore}): deben explorarse todas.}
\end{columns}
\end{frame}

%=======================================================================
% FILTRADO DE ERRORES
%=======================================================================
\begin{frame}{Filtrado -- Transiciones controlables a error}
\begin{columns}[c]
\column{0.44\textwidth}
\begin{center}
\begin{tikzpicture}[node distance=1.9cm]
  \node[st, fill=gray!12] (s)  {$s$};
  \node[st, above right=0.8cm and 1.8cm of s] (t1) {$t_1$};
  \node[errst, below right=0.8cm and 1.8cm of s] (e) {\footnotesize ERR};
  \draw[tChosen] (s) -- node[above, lbl, sloped]{$a$} (t1);
  \draw[tError]  (s) -- node[below, lbl, sloped]{$b$} (e);
  \node[below=0.05cm of s, cChosen, font=\tiny] {actual};
\end{tikzpicture}
\end{center}

\column{0.52\textwidth}
\textbf{Transiciones controlables a error} son descartadas antes de aplicar cualquier regla.

\medskip
\begin{itemize}
  \item[\textcolor{cChosen}{$\blacktriangleright$}] Elegir $a \to t_1$
  \item[\textcolor{cError}{$\times$}] $b \to \text{ERR}$ descartada
\end{itemize}

\medskip
{\small Un estado es \emph{error} si est\'a en el conjunto de errores del contexto, o si alg\'un sub-estado del producto paralelo es \texttt{``ERROR''}.}

\medskip
{\small Las no-controlables a error \textbf{no} se filtran (el director debe expandirlas).}
\end{columns}
\end{frame}

%=======================================================================
% FALLBACK
%=======================================================================
\begin{frame}{Fallback -- Pila e Ignorados}

Cuando no hay transiciones desde el \textbf{estado actual} en la \emph{frontera}:

\bigskip
\begin{enumerate}
  \item[\textcolor{cStack}{\textbf{1.}}] \textcolor{cStack}{\textbf{Pop stack}} -- tomar la transici\'on del tope de la pila.\\
    {\small Pila $\to$ no-controlables diferidas de pasos anteriores.}

  \bigskip
  \item[\textcolor{cIgnore}{\textbf{2.}}] \textcolor{cIgnore}{\textbf{Pop ignore}} -- si la pila est\'a vac\'ia, tomar de la lista de ignorados.\\
    {\small Ignorados $\to$ controlables descartadas en decisiones anteriores.}

  \bigskip
  \item[\textbf{3.}] \textbf{Primera de la frontera} -- \'ultimo recurso si ambas listas est\'an vac\'ias.
\end{enumerate}

\bigskip
\begin{alertblock}{Invariante}
Pila $\to$ ignorados garantiza que todas las ramas no-controlables se exploran \emph{antes} de revisar las decisiones controlables diferidas.
\end{alertblock}
\end{frame}

%=======================================================================
% ESTADOS Y SUB-ESTADOS
%=======================================================================
\begin{frame}{Estados de la Planta}

\textbf{Estado de la planta:} tupla de sub-estados, uno por cada LTS componente.
\[
  s = (s_0,\; s_1,\; \ldots,\; s_{m-1})
\]

\bigskip
\textbf{Ejemplo:}
\[
  \texttt{Wait[2] | Wait[4] | Next[1] | Next[0] | off}
\]

\medskip
Cada elemento es un sub-estado: \texttt{Wait[2]} es el sub-estado del primer LTS, \texttt{Wait[4]} el del segundo, etc.

\bigskip
\begin{alertblock}{Información valiosa}
Los sub-estados son strings con semántica del componente. El \texttt{picker} los lee para decidir qué transición controlable expandir.
\end{alertblock}

\end{frame}

%=======================================================================
% PICKERS POR FAMILIA
%=======================================================================
\begin{frame}{}
\begin{center}
\Huge \texttt{pickControllable} por familia
\end{center}
\end{frame}

%=======================================================================
% PICKER: DP
%=======================================================================
\begin{frame}{\texttt{pickControllable} -- DP (Dining Philosophers)}

\textbf{Candidatas}: \texttt{take[i][j]} (fil\'osofo $i$ toma el fork $j$).

\bigskip
\textbf{Prioridad}:
\begin{enumerate}
  \item \textbf{Ready}: elegir \texttt{take[i]} con menor $i$ tal que el fil\'osofo $i$ est\'a en \texttt{``Ready''}.
  \item \textbf{Hungry}: elegir \texttt{take[i]} con menor $i$ tal que el fil\'osofo $i$ est\'a en \texttt{``Hungry''}\\
    y el monitor \emph{no} est\'a en \texttt{``Done''}.
  \item \textbf{Fallback}: primera transici\'on.
\end{enumerate}

\bigskip
\begin{block}{Intuici\'on}
\texttt{Ready} significa que el fil\'osofo ya obtuvo el primer fork y puede comer si toma el segundo. Atenderlo primero minimiza la espera. El \'indice m\'inimo da orden determin\'istico.
\end{block}
\end{frame}

%=======================================================================
% PICKER: TL
%=======================================================================
\begin{frame}{\texttt{pickControllable} -- TL (Transfer Line)}

\textbf{Candidatas}: \texttt{get[i]} (recoger pieza en posici\'on $i$).

\bigskip
\textbf{Prioridad}:
\begin{enumerate}
  \item \textbf{Ya explorado}: elegir \texttt{get[i]} cuyo estado destino ya fue visitado.\\
    {\small Cierra el ciclo de la l\'inea de transferencia.}
  \item \textbf{Mayor \'indice}: elegir \texttt{get[i]} con $i$ m\'aximo (que no vaya a error).\\
    {\small Avanza lo m\'as posible hacia el final de la cadena.}
  \item \textbf{Fallback}: primera transici\'on.
\end{enumerate}

\bigskip
\begin{block}{Intuici\'on}
La l\'inea tiene un orden lineal natural. Cerrar el ciclo primero; si no es posible, avanzar al m\'aximo $i$ sin error.
\end{block}
\end{frame}

%=======================================================================
% PICKER: AT
%=======================================================================
\begin{frame}{\texttt{pickControllable} -- AT (Air Traffic)}

\textbf{Candidatas}: \texttt{approach[i]}, \texttt{descend[i][j]}.

\bigskip
\textbf{Prioridad}:
\begin{enumerate}
  \item \textbf{Approach}: elegir \texttt{approach[i]} si las alturas del estado actual son \emph{consecutivas}\\
    (sin huecos vac\'ios entre aviones).
  \item \textbf{Descend}: elegir \texttt{descend[i][j]} con menor $j$ tal que:
    \begin{itemize}
      \item el avi\'on $i$ est\'a en altura $j{+}1$ (justo encima del objetivo), y
      \item la altura $j$ est\'a vac\'ia (\texttt{``Empty''}) en el estado actual.
    \end{itemize}
  \item \textbf{Fallback}: primera transici\'on.
\end{enumerate}

\bigskip
\begin{block}{Intuici\'on}
Primero completar aproximaciones cuando los aviones est\'an ordenados; luego bajar el avi\'on que est\'a justo encima del primer hueco libre.
\end{block}
\end{frame}

%=======================================================================
% PICKER: CM
%=======================================================================
\begin{frame}{\texttt{pickControllable} -- CM (Cat and Mouse)}

\textbf{Candidatas}: \texttt{mouse[i][move[j]]} (rat\'on $i$ se mueve a posici\'on $j$).

\bigskip
\textbf{Regla}:
\[
  \text{safePlace} = \left\lfloor\frac{2k+1}{2}\right\rfloor
\]
\begin{itemize}
  \item Elegir el movimiento que minimiza $|j - \text{safePlace}|$ (destino m\'as cercano al centro).
  \item Empate: menor \'indice de rat\'on $i$.
  \item \textbf{Fallback}: primera transici\'on.
\end{itemize}

\bigskip
\begin{block}{Intuici\'on}
El rat\'on debe alcanzar el centro del tablero (posici\'on segura / goal). Mover siempre hacia la posici\'on central minimiza los pasos hasta el objetivo.
\end{block}
\end{frame}

%=======================================================================
% PICKER: BW
%=======================================================================
\begin{frame}{\texttt{pickControllable} -- BW (Bidding Workflow)}

\textbf{Candidatas}: \texttt{approve}, \texttt{refuse}, \texttt{assign[i]}.

\bigskip
\textbf{Prioridad}:
\begin{enumerate}
  \item \textbf{Approve}: elegir \texttt{approve} si no lleva a error.
  \item \textbf{Refuse}: elegir \texttt{refuse} si el primer sub-estado del estado \emph{origen} es \texttt{``Rejected''}.
  \item \textbf{Assign}: elegir \texttt{assign[i]} con menor $i$ tal que:
    \begin{itemize}
      \item el crew $i$ est\'a en \texttt{``Pending''} o \texttt{``Rejected[x]''} en el estado actual, y
      \item no lleva a error.
    \end{itemize}
  \item \textbf{Fallback}: primera transici\'on.
\end{enumerate}

\bigskip
\begin{block}{Intuici\'on}
Aprobar y rechazar correctamente antes de asignar; dentro de las asignaciones, el orden por \'indice da determinismo.
\end{block}
\end{frame}

%=======================================================================
% PICKER: TA (placeholder)
%=======================================================================
\begin{frame}{\texttt{pickControllable} -- TA (Travel Agency)}

\textbf{Candidatas}: \texttt{purchase[i]}, \texttt{cancel[i]}, \texttt{agency.succ}, \texttt{agency.fail}.

{\small Layout: \texttt{[Agency | AgencyMonitor | Svc(0) | SvcMon(0) | \ldots | Svc(n{-}1) | SvcMon(n{-}1) | fluent]}}

\medskip
\textbf{Prioridad}:
\begin{enumerate}
  \item \textbf{All Success}: si todos los sub-estados \texttt{SvcMon(i)} son \texttt{``Success''} y \texttt{agency.succ} no lleva a error $\to$ expandir \texttt{agency.succ}.
  \item \textbf{agency.fail}: si \texttt{agency.fail} no lleva a error $\to$ expandirla.
  \item \textbf{purchase[i]}: elegir \texttt{purchase[i]} con $i$ m\'as cercano al \'indice $j$ del sub-estado \texttt{AgencyMonitor = ``Disallow[j]''}, tal que \texttt{SvcMon(i)} $\ne$ \texttt{``Success''}.
  \item \textbf{Fallback}: primera transici\'on.
\end{enumerate}

\medskip
\begin{block}{Intuici\'on}
Confirmar la compra cuando todos los servicios aceptaron; si no, comprar el servicio m\'as cercano al \'indice actualmente denegado por el monitor de la agencia.
\end{block}
\end{frame}

%=======================================================================
% HEATMAP MACRO
% Cell (n,k): x=k-1..k, y=n-1..n  (n=1 at bottom, n=15 at top)
% #1 = title, #2 = fill commands
%=======================================================================
\newcommand{\cmap}[2]{%
  \begin{tikzpicture}[scale=0.135]
    \node[above, font=\bfseries\scriptsize, inner sep=1.5pt] at (7.5,15) {#1};
    \fill[black] (0,0) rectangle (15,15);
    #2
    \draw[gray!30, ultra thin, step=1] (0,0) grid (15,15);
    \draw[black, thin] (0,0) rectangle (15,15);
    \node[font=\tiny, below] at (0.5,-0.3)  {\tiny 1};
    \node[font=\tiny, below] at (14.5,-0.3) {\tiny 15};
    \node[font=\tiny, left]  at (-0.3,0.5)  {\tiny 1};
    \node[font=\tiny, left]  at (-0.3,14.5) {\tiny 15};
    \node[font=\tiny, below] at (7.5,-1.2) {\tiny$k$};
    \node[font=\tiny, left, rotate=90] at (-1.5,7.5) {\tiny$n$};
  \end{tikzpicture}%
}

% Legend helper
\newcommand{\hmlegend}{%
  \vspace{2pt}
  {\tiny
    \textcolor{cB1}{$\blacksquare$}~$\le100$\quad
    \textcolor{cB2}{$\blacksquare$}~$\le999$\quad
    \textcolor{cB3}{$\blacksquare$}~$\le5000$\quad
    \textcolor{cB4}{$\blacksquare$}~$\le15000$\quad
    \textcolor{white}{$\blacksquare$}\llap{\colorbox{black}{\phantom{X}}}~timeout
  }
}

%=======================================================================
% RESULTADOS
%=======================================================================
\begin{frame}{}
\begin{center}
\Huge Resultados comparativos\\[0.5em]
\large Budget: 15\,000 expansiones por instancia
\end{center}
\end{frame}

%=======================================================================
% RA heatmaps (6 families, 3+3 layout)
%=======================================================================
\begin{frame}{RA -- Instancias resueltas (budget = 15\,000)}
\begin{center}
\vspace{-6pt}
\cmap{AT (81/225)}{%
  \fill[cB1] (0,0) rectangle (15,2);
  \fill[cB1] (0,2) rectangle (1,3);
  \fill[cB2] (1,2) rectangle (15,3);
  \fill[cB1] (0,3) rectangle (1,4);
  \fill[cB2] (1,3) rectangle (2,4);
  \fill[cB3] (2,3) rectangle (3,4);
  \fill[cB2] (3,3) rectangle (15,4);
  \fill[cB2] (0,4) rectangle (1,5);
  \fill[cB3] (1,4) rectangle (2,5);
  \fill[cB4] (2,4) rectangle (4,5);
  \fill[cB3] (4,4) rectangle (15,5);
  \fill[cB2] (0,5) rectangle (1,6);
  \fill[cB4] (1,5) rectangle (2,6);
  \fill[cB3] (0,6) rectangle (1,7);
  \fill[cB3] (0,7) rectangle (1,8);
  \fill[cB4] (0,8) rectangle (1,9);
  \fill[cB4] (0,9) rectangle (1,10);
}%
\hspace{6pt}%
\cmap{BW (146/225)}{%
  \fill[cB1] (0,0) rectangle (15,2);
  \fill[cB1] (0,2) rectangle (6,3);
  \fill[cB2] (6,2) rectangle (15,3);
  \fill[cB1] (0,3) rectangle (2,4);
  \fill[cB2] (2,3) rectangle (15,4);
  \fill[cB2] (0,4) rectangle (15,5);
  \fill[cB2] (0,5) rectangle (8,6);
  \fill[cB3] (8,5) rectangle (15,6);
  \fill[cB2] (0,6) rectangle (3,7);
  \fill[cB3] (3,6) rectangle (15,7);
  \fill[cB3] (0,7) rectangle (11,8);
  \fill[cB4] (11,7) rectangle (15,8);
  \fill[cB3] (0,8) rectangle (4,9);
  \fill[cB4] (4,8) rectangle (15,9);
  \fill[cB3] (0,9) rectangle (1,10);
  \fill[cB4] (1,9) rectangle (8,10);
  \fill[cB4] (0,10) rectangle (3,11);
}%
\hspace{6pt}%
\cmap{CM (23/225)}{%
  \fill[cB1] (0,0) rectangle (3,1);
  \fill[cB2] (3,0) rectangle (15,1);
  \fill[cB2] (0,1) rectangle (3,2);
  \fill[cB3] (3,1) rectangle (4,2);
  \fill[cB2] (0,2) rectangle (1,3);
  \fill[cB3] (1,2) rectangle (2,3);
  \fill[cB4] (2,2) rectangle (3,3);
  \fill[cB4] (0,3) rectangle (1,4);
}%
\\[3pt]
\cmap{DP (165/225)}{%
  \fill[cB1] (0,0) rectangle (15,3);
  \fill[cB1] (0,3) rectangle (8,4);
  \fill[cB2] (8,3) rectangle (15,4);
  \fill[cB2] (0,4) rectangle (15,7);
  \fill[cB3] (0,7) rectangle (15,9);
  \fill[cB4] (0,9) rectangle (15,11);
}%
\hspace{6pt}%
\cmap{TL (225/225)}{%
  \fill[cB1] (0,0) rectangle (15,15);
  \fill[cB2] (0,5) rectangle (1,13);
  \fill[cB3] (0,13) rectangle (1,15);
}%
\hspace{6pt}%
\cmap{TA (60/225)}{%
  \fill[cB1] (0,0) rectangle (15,1);
  \fill[cB2] (0,1) rectangle (15,3);
  \fill[cB3] (0,3) rectangle (15,4);
}%

\vspace{3pt}
\hmlegend
\end{center}
\end{frame}

%=======================================================================
% SuperDFS heatmaps (6 families, 3+3 layout)
%=======================================================================
\begin{frame}{SuperDFS -- Instancias resueltas (budget = 15\,000)}
\begin{center}
\vspace{-6pt}
\cmap{AT (85/225)}{%
  \fill[cB1] (0,0) rectangle (15,4);
  \fill[cB1] (0,4) rectangle (1,5);
  \fill[cB3] (1,4) rectangle (2,5);
  \fill[cB4] (2,4) rectangle (3,5);
  \fill[cB3] (4,4) rectangle (15,5);
  \fill[cB1] (0,5) rectangle (1,6);
  \fill[cB4] (1,5) rectangle (2,6);
  \fill[cB1] (0,6) rectangle (1,15);
}%
\hspace{6pt}%
\cmap{BW (146/225)}{%
  \fill[cB1] (0,0) rectangle (15,2);
  \fill[cB1] (0,2) rectangle (6,3);
  \fill[cB2] (6,2) rectangle (15,3);
  \fill[cB1] (0,3) rectangle (2,4);
  \fill[cB2] (2,3) rectangle (15,4);
  \fill[cB2] (0,4) rectangle (15,5);
  \fill[cB2] (0,5) rectangle (8,6);
  \fill[cB3] (8,5) rectangle (15,6);
  \fill[cB2] (0,6) rectangle (3,7);
  \fill[cB3] (3,6) rectangle (15,7);
  \fill[cB3] (0,7) rectangle (11,8);
  \fill[cB4] (11,7) rectangle (15,8);
  \fill[cB3] (0,8) rectangle (4,9);
  \fill[cB4] (4,8) rectangle (15,9);
  \fill[cB3] (0,9) rectangle (1,10);
  \fill[cB4] (1,9) rectangle (8,10);
  \fill[cB4] (0,10) rectangle (3,11);
}%
\hspace{6pt}%
\cmap{CM (30/225)}{%
  \fill[cB1] (0,0) rectangle (3,1);
  \fill[cB2] (3,0) rectangle (15,1);
  \fill[cB2] (0,1) rectangle (3,2);
  \fill[cB3] (3,1) rectangle (6,2);
  \fill[cB4] (6,1) rectangle (11,2);
  \fill[cB2] (0,2) rectangle (1,3);
  \fill[cB3] (1,2) rectangle (2,3);
  \fill[cB4] (2,2) rectangle (3,3);
  \fill[cB4] (0,3) rectangle (1,4);
}%
\\[3pt]
\cmap{DP (165/225)}{%
  \fill[cB1] (0,0) rectangle (15,3);
  \fill[cB1] (0,3) rectangle (9,4);
  \fill[cB2] (9,3) rectangle (15,4);
  \fill[cB2] (0,4) rectangle (15,8);
  \fill[cB3] (0,8) rectangle (15,10);
  \fill[cB4] (0,10) rectangle (15,11);
}%
\hspace{6pt}%
\cmap{TL (225/225)}{%
  \fill[cB1] (0,0) rectangle (15,15);
}%
\hspace{6pt}%
\cmap{TA (129/225)}{%
  \fill[cB1] (0,0) rectangle (15,1);
  \fill[cB1] (0,1) rectangle (6,2);
  \fill[cB2] (6,1) rectangle (15,2);
  \fill[cB2] (0,2) rectangle (15,4);
  \fill[cB2] (0,4) rectangle (7,5);
  \fill[cB3] (7,4) rectangle (15,5);
  \fill[cB2] (0,5) rectangle (2,6);
  \fill[cB3] (2,5) rectangle (15,6);
  \fill[cB3] (0,6) rectangle (10,7);
  \fill[cB4] (10,6) rectangle (15,7);
  \fill[cB3] (0,7) rectangle (3,8);
  \fill[cB4] (3,7) rectangle (15,8);
  \fill[cB4] (0,8) rectangle (7,9);
  \fill[cB4] (0,9) rectangle (2,10);
}%

\vspace{3pt}
\hmlegend
\end{center}
\end{frame}

%=======================================================================
% GRAFICO COMPARATIVO
%=======================================================================
\begin{frame}[shrink=2]{Comparativa -- Instancias resueltas (budget = 15\,000)}
\begin{center}
\begin{tikzpicture}
\begin{axis}[
    ybar stacked,
    bar width=1.8cm,
    symbolic x coords={RA, SuperDFS},
    xtick=data,
    x tick label style={font=\large\bfseries},
    ylabel={Instancias resueltas},
    ylabel style={font=\normalsize},
    ymin=0, ymax=900,
    width=9cm, height=7.1cm,
    legend style={at={(1.02,1)}, anchor=north west, font=\scriptsize},
    legend cell align=left,
    enlarge x limits=0.5,
    xtick style={draw=none},
    ymajorgrids=true,
    grid style={dashed, gray!40},
    nodes near coords,
    nodes near coords align={center},
    nodes near coords style={font=\scriptsize\bfseries, text=white},
    every node near coord/.append style={yshift=0pt},
]
\addplot+[fill=blue!50,   draw=blue!80!black]    coordinates {(RA,81)  (SuperDFS,85)};
\addlegendentry{AT}
\addplot+[fill=orange!70, draw=orange!90!black]  coordinates {(RA,146) (SuperDFS,146)};
\addlegendentry{BW}
\addplot+[fill=red!60,    draw=red!80!black]     coordinates {(RA,23)  (SuperDFS,30)};
\addlegendentry{CM}
\addplot+[fill=green!60!black, draw=green!80!black] coordinates {(RA,165) (SuperDFS,165)};
\addlegendentry{DP}
\addplot+[fill=purple!60, draw=purple!80!black]  coordinates {(RA,225) (SuperDFS,225)};
\addlegendentry{TL}
\addplot+[fill=teal!70,   draw=teal!90!black]    coordinates {(RA,60)  (SuperDFS,129)};
\addlegendentry{TA}
\node[above, font=\normalsize\bfseries] at (axis cs:RA, 700) {700};
\node[above, font=\normalsize\bfseries] at (axis cs:SuperDFS, 780) {780};
\end{axis}
\end{tikzpicture}
\end{center}
\end{frame}

%=======================================================================
% SUPER RL -- SEPARADOR
%=======================================================================
\begin{frame}{}
\begin{center}
\Huge Super RL
\end{center}
\end{frame}

%=======================================================================
% SUPER RL -- QUÉ ES
%=======================================================================
\begin{frame}{Super RL -- Idea}

\begin{block}{Objetivo}
Reemplazar el \texttt{picker} manual de SuperDFS por una red neuronal entrenada por RL.
\end{block}

\bigskip
\textbf{Arquitectura:}
\begin{itemize}
  \item Red MLP: recibe el \textbf{vector de features} de cada transición candidata y devuelve un score.
  \item En cada paso se elige la transición con mayor score (\emph{argmax}).
  \item Una red por familia (AT, DP, BW, CM, TL, TA).
\end{itemize}

\bigskip
\begin{alertblock}{Mismo esqueleto que SuperDFS}
El algoritmo de exploración (pila, ignorados, fallback) es idéntico. Solo cambia el \texttt{picker}: en lugar de reglas codificadas a mano, usa la red.
\end{alertblock}

\end{frame}

%=======================================================================
% SUPER_CUSTOM FEATURES
%=======================================================================
\begin{frame}{Features: \texttt{SUPER\_CUSTOM}}

\begin{center}
\renewcommand{\arraystretch}{1.25}
\begin{tabular}{c|c|l}
\toprule
Familia & \# & Features \\
\midrule
AT & 5 & \texttt{is\_approach, target\_empty, consecutive,} \\
   &   & \texttt{height\_norm, height\_matches\_slot} \\
\midrule
DP & 5 & \texttt{is\_take, ready, hungry, min\_ready\_idx, min\_hungry\_idx} \\
\midrule
BW & 7 & \texttt{is\_approve, is\_refuse, is\_assign, doc\_rejected,} \\
   &   & \texttt{crew\_pending, crew\_rejected, min\_eligible} \\
\midrule
CM & 3 & \texttt{dist\_to\_safe\_norm, is\_min\_dist, is\_min\_mouse} \\
\midrule
TL & 2 & \texttt{dest\_explored, is\_max\_idx} \\
\midrule
TA & 6 & \texttt{is\_succ, is\_fail, is\_purchase,} \\
   &   & \texttt{all\_success, this\_success, dist\_to\_disallow\_norm} \\
\bottomrule
\end{tabular}
\end{center}

\end{frame}

%=======================================================================
% ENTRENAMIENTO -- PathLearning
%=======================================================================
\begin{frame}{Entrenamiento: PathLearning}

\textbf{Objetivo:} aprender a imitar el mejor camino encontrado.

\bigskip
\begin{enumerate}
  \item \textbf{Ronda 0 -- exploración aleatoria:} correr $N$ episodios aleatorios. Guardar el episodio con mayor reward (= menos expansiones). \textbf{Ese es el mejor path.}

  \bigskip
  \item \textbf{Rondas $k > 0$ -- guided:} correr $N$ episodios con la red actual ($\varepsilon$-greedy). Si alguno supera el mejor path global, actualizarlo.

  \bigskip
  \item \textbf{Entrenamiento por ronda:} construir un dataset desde el mejor path y re-entrenar la red desde cero hasta convergencia.
\end{enumerate}

\end{frame}

%=======================================================================
% ENTRENAMIENTO -- Ranking de feature vectors
%=======================================================================
\begin{frame}{Entrenamiento: Ranking de Feature Vectors}

\textbf{¿Cómo construir el dataset desde el mejor path?}

\bigskip
Para cada feature vector (FV) único observado, se rastrean:
\begin{itemize}
  \item $\text{gano}$: pasos donde ese FV fue \emph{elegido}
  \item $\text{perdio}$: pasos donde ese FV estaba pero \emph{no fue elegido}
\end{itemize}

\medskip
\textbf{Dominancia:} $A$ domina a $B$ si $A$ ganó en algún paso donde $B$ estaba y perdió.

\medskip
\textbf{Target:}
\[
  y = \begin{cases}
    -1 & \text{si } \text{gano} = \emptyset \\
    |\text{gano}| - |\{B : B \text{ domina a } A\}| & \text{si no}
  \end{cases}
\]

\medskip
Los targets se colapsan a rangos densos. La red se entrena con MSELoss; en inferencia se usa \emph{argmax}.

\end{frame}

%=======================================================================
% GRAFICO COMPARATIVO 3 MODELOS
%=======================================================================
\begin{frame}[shrink=2]{Comparativa -- RA vs SuperDFS vs SuperRL (budget = 15\,000)}
\begin{center}
\begin{tikzpicture}
\begin{axis}[
    ybar stacked,
    bar width=1.3cm,
    symbolic x coords={RA, SuperDFS, SuperRL},
    xtick=data,
    x tick label style={font=\large\bfseries},
    ylabel={Instancias resueltas},
    ylabel style={font=\normalsize},
    ymin=0, ymax=900,
    width=10cm, height=7.1cm,
    legend style={at={(1.02,1)}, anchor=north west, font=\scriptsize},
    legend cell align=left,
    enlarge x limits=0.4,
    xtick style={draw=none},
    ymajorgrids=true,
    grid style={dashed, gray!40},
    nodes near coords,
    nodes near coords align={center},
    nodes near coords style={font=\scriptsize\bfseries, text=white},
    every node near coord/.append style={yshift=0pt},
]
\addplot+[fill=blue!50,   draw=blue!80!black]    coordinates {(RA,81)  (SuperDFS,85)  (SuperRL,85)};
\addlegendentry{AT}
\addplot+[fill=orange!70, draw=orange!90!black]  coordinates {(RA,146) (SuperDFS,146) (SuperRL,165)};
\addlegendentry{BW}
\addplot+[fill=red!60,    draw=red!80!black]     coordinates {(RA,23)  (SuperDFS,30)  (SuperRL,30)};
\addlegendentry{CM}
\addplot+[fill=green!60!black, draw=green!80!black] coordinates {(RA,165) (SuperDFS,165) (SuperRL,165)};
\addlegendentry{DP}
\addplot+[fill=purple!60, draw=purple!80!black]  coordinates {(RA,225) (SuperDFS,225) (SuperRL,225)};
\addlegendentry{TL}
\addplot+[fill=teal!70,   draw=teal!90!black]    coordinates {(RA,60)  (SuperDFS,129) (SuperRL,129)};
\addlegendentry{TA}
\node[above, font=\normalsize\bfseries] at (axis cs:RA, 700)      {700};
\node[above, font=\normalsize\bfseries] at (axis cs:SuperDFS, 780) {780};
\node[above, font=\normalsize\bfseries] at (axis cs:SuperRL, 799)  {799};
\end{axis}
\end{tikzpicture}
\end{center}
\end{frame}

%=======================================================================
% SUPER RL -- HEATMAPS
%=======================================================================
\begin{frame}{SuperRL -- Instancias resueltas (budget = 15\,000)}
\begin{center}
\vspace{-6pt}
\cmap{AT (85/225)}{%
  \fill[cB1] (0,0) rectangle (15,2);
  \fill[cB1] (0,2) rectangle (1,3);
  \fill[cB2] (1,2) rectangle (15,3);
  \fill[cB1] (0,3) rectangle (1,4);
  \fill[cB2] (1,3) rectangle (2,4);
  \fill[cB3] (2,3) rectangle (3,4);
  \fill[cB2] (3,3) rectangle (7,4);
  \fill[cB3] (7,3) rectangle (15,4);
  \fill[cB1] (0,4) rectangle (1,5);
  \fill[cB3] (1,4) rectangle (2,5);
  \fill[cB4] (2,4) rectangle (3,5);
  \fill[cB4] (4,4) rectangle (15,5);
  \fill[cB1] (0,5) rectangle (1,6);
  \fill[cB4] (1,5) rectangle (2,6);
  \fill[cB1] (0,6) rectangle (1,15);
}%
\hspace{6pt}%
\cmap{BW (165/225)}{%
  \fill[cB1] (0,0) rectangle (15,2);
  \fill[cB1] (0,2) rectangle (8,3);
  \fill[cB2] (8,2) rectangle (15,3);
  \fill[cB1] (0,3) rectangle (4,4);
  \fill[cB2] (4,3) rectangle (15,4);
  \fill[cB2] (0,4) rectangle (15,7);
  \fill[cB3] (0,7) rectangle (15,10);
  \fill[cB4] (0,10) rectangle (15,11);
}%
\hspace{6pt}%
\cmap{CM (30/225)}{%
  \fill[cB1] (0,0) rectangle (3,1);
  \fill[cB2] (3,0) rectangle (15,1);
  \fill[cB2] (0,1) rectangle (3,2);
  \fill[cB3] (3,1) rectangle (6,2);
  \fill[cB4] (6,1) rectangle (11,2);
  \fill[cB2] (0,2) rectangle (1,3);
  \fill[cB3] (1,2) rectangle (2,3);
  \fill[cB4] (2,2) rectangle (3,3);
  \fill[cB4] (0,3) rectangle (1,4);
}%
\\[3pt]
\cmap{DP (165/225)}{%
  \fill[cB1] (0,0) rectangle (15,3);
  \fill[cB1] (0,3) rectangle (10,4);
  \fill[cB2] (10,3) rectangle (15,4);
  \fill[cB2] (0,4) rectangle (15,8);
  \fill[cB3] (0,8) rectangle (15,10);
  \fill[cB4] (0,10) rectangle (15,11);
}%
\hspace{6pt}%
\cmap{TL (225/225)}{%
  \fill[cB1] (0,0) rectangle (15,15);
}%
\hspace{6pt}%
\cmap{TA (129/225)}{%
  \fill[cB1] (0,0) rectangle (15,1);
  \fill[cB1] (0,1) rectangle (6,2);
  \fill[cB2] (6,1) rectangle (15,2);
  \fill[cB2] (0,2) rectangle (15,4);
  \fill[cB2] (0,4) rectangle (7,5);
  \fill[cB3] (7,4) rectangle (15,5);
  \fill[cB2] (0,5) rectangle (2,6);
  \fill[cB3] (2,5) rectangle (15,6);
  \fill[cB3] (0,6) rectangle (10,7);
  \fill[cB4] (10,6) rectangle (15,7);
  \fill[cB3] (0,7) rectangle (3,8);
  \fill[cB4] (3,7) rectangle (15,8);
  \fill[cB4] (0,8) rectangle (7,9);
  \fill[cB4] (0,9) rectangle (2,10);
}%

\vspace{3pt}
\hmlegend
\end{center}
\end{frame}

\end{document}
